\textbf{\textcolor{CExample}{例 1（基础）}}
\n\n\textbf{\textcolor{CExample}{题目：}} 已知三棱锥 $P$-$ABC$ 的体积 $V=30$，底面 $\triangle ABC$ 的面积 $S=18$，求顶点 $P$ 到底面 $ABC$ 的距离 $h$。\n\n\textbf{\textcolor{CExample}{解题步骤：}}\n
\begin{description}[style=nextline, leftmargin=2.8em, labelsep=0.5em, itemsep=0.45em]
	\n
	\item[\textbf{第一步（列条件已知量）：}]
		体积 $V=30$，底面面积 $S=18$，待求距离 $h$。\n
		\par\textbf{理由：}\n 等体积法公式 $h=\dfrac{3V}{S}$ 需要这两个量。\n
	\item[\textbf{第二步（代入公式）：}]
		\n
		\[
			h = \frac{3 \times 30}{18}.
		\]
		\n
		\par\textbf{理由：}\n 由等体积法公式 $h=\dfrac{3V}{S}$ 直接代入。\n
	\item[\textbf{第三步（计算结果）：}]
		\n
		\[
			\begin{aligned}
				h &= \frac{90}{18} \\ &= 5.
		\end{aligned}
		\]
		\n
		\par\textbf{理由：}\n 算数运算。\n
\end{description}
\n\n\textbf{\textcolor{CExample}{关键结论：}} \textbf{点 $P$ 到平面 $ABC$ 的距离为 $5$。}\n\n

\medskip
\n\textbf{\textcolor{CExample}{例 2（稍变形）}}\n\n\textbf{\textcolor{CExample}{题目：}} 三棱锥 $A$-$BCD$ 中，已知底面 $BCD$ 面积 $S_1=12$，顶点 $A$ 到底面 $BCD$ 的距离 $h_1=4$；若 $\triangle ACD$ 面积 $S_2=8$，求顶点 $B$ 到底面 $ACD$ 的距离 $h_2$。\n\n\textbf{\textcolor{CExample}{解题步骤：}}\n
\begin{description}[style=nextline, leftmargin=2.8em, labelsep=0.5em, itemsep=0.45em]
	\n
	\item[\textbf{第一步（求三棱锥体积）：}]
		以 $BCD$ 为底面，$A$ 为顶点，则\n
		\[
			\begin{aligned}
				V &= \frac13 S_1 h_1 \\ &= \frac13 \times 12 \times 4 \\ &= 16.
		\end{aligned}
		\]
		\n
		\par\textbf{理由：}\n 三棱锥体积公式 $V=\dfrac13 S h$。\n
	\item[\textbf{第二步（换底建立等式）：}]
		对于同一个三棱锥，若以 $\triangle ACD$ 为底面，$B$ 为顶点，体积仍为 $V$，设 $h_2$ 为 $B$ 到平面 $ACD$ 的距离，则\n
		\[
			V = \frac13 S_2 h_2.
		\]
		\n
		\par\textbf{理由：}\n 等体积法：同一个三棱锥体积不变，更换底面后对应的高改变。\n
	\item[\textbf{第三步（解出 $h_2$）：}]
		\n
		\[
			\begin{aligned}
				h_2 &= \frac{3V}{S_2} \\ &= \frac{3 \times 16}{8} \\ &= 6.
		\end{aligned}
		\]
		\n
		\par\textbf{理由：}\n 由 $V=\dfrac13 S_2 h_2$ 解出 $h_2$。\n
\end{description}
\n\n\textbf{\textcolor{CExample}{关键结论：}} \textbf{点 $B$ 到平面 $ACD$ 的距离为 $6$。}\n\n

\medskip
\n\textbf{\textcolor{CExample}{例 3（综合应用）}}\n\n\textbf{\textcolor{CExample}{题目：}} 三棱锥 $P$-$ABC$ 中，$PA,\ PB,\ PC$ 两两垂直，且 $PA=PB=PC=2$，求点 $P$ 到平面 $ABC$ 的距离。\n\n\textbf{\textcolor{CExample}{解题步骤：}}\n
\begin{description}[style=nextline, leftmargin=2.8em, labelsep=0.5em, itemsep=0.45em]
	\n
	\item[\textbf{第一步（求三棱锥体积）：}]
		以 $\triangle PAB$ 为底面，$PC$ 为高（因为 $PC \perp$ 平面 $PAB$），则\n
		\[
			\begin{aligned}
				V &= \frac13 \times \frac12 \times PA \times PB \times PC \\ &= \frac13 \times \frac12 \times 2 \times 2 \times 2 \\ &= \frac{4}{3}.
		\end{aligned}
		\]
		\n
		\par\textbf{理由：}\n $PA,\ PB,\ PC$ 两两垂直，故 $PC$ 垂直于底面 $PAB$，底面 $PAB$ 为直角三角形，面积 $\dfrac12 \cdot PA \cdot PB$。\n
	\item[\textbf{第二步（求底面 $\triangle ABC$ 的面积）：}]
		由勾股定理，\n
		\[
			\begin{aligned}
				AB &= \sqrt{PA^2+PB^2} \\ &= 2\sqrt{2},\\ BC &= \sqrt{PB^2+PC^2} \\ &= 2\sqrt{2},\\ AC &= \sqrt{PA^2+PC^2} \\ &= 2\sqrt{2},
		\end{aligned}
		\]
		\n 所以 $\triangle ABC$ 是边长为 $2\sqrt{2}$ 的等边三角形，\n
		\[
			\begin{aligned}
				S_{\triangle ABC} &= \frac{\sqrt3}{4} \times (2\sqrt{2})^2 \\ &= \frac{\sqrt3}{4} \times 8 \\ &= 2\sqrt{3}.
		\end{aligned}
		\]
		\n
		\par\textbf{理由：}\n 垂直关系导出各棱长，等边三角形面积公式。\n
	\item[\textbf{第三步（等体积法求距离）：}]
		设 $h$ 为点 $P$ 到平面 $ABC$ 的距离，则由等体积法\n
		\[
			\begin{aligned}
				h &= \frac{3V}{S_{\triangle ABC}} \\ &= \frac{3 \times \dfrac{4}{3}}{2\sqrt{3}} \\ &= \frac{4}{2\sqrt{3}} \\ &= \frac{2}{\sqrt{3}} \\ &= \frac{2\sqrt{3}}{3}.
		\end{aligned}
		\]
		\n
		\par\textbf{理由：}\n 等体积法公式 $h=\dfrac{3V}{S}$，将已求得的 $V$ 和 $S_{\triangle ABC}$ 代入。\n
\end{description}
\n\n\textbf{\textcolor{CExample}{关键结论：}} \textbf{点 $P$ 到平面 $ABC$ 的距离为 $\dfrac{2\sqrt{3}}{3}$。}
